\chapter{Documentación y Presentación}\label{cap:documentacion}

\section{Elaboración de la documentación técnica}

La tercera fase del proyecto se ha dedicado íntegramente a la consolidación de toda
la información generada durante el ciclo de vida del desarrollo para cumplir con los
requisitos de evaluación de la asignatura. El núcleo de esta fase es la redacción de
la presente memoria, un documento que no solo describe el producto final, sino que
justifica cada decisión tomada por el equipo, desde la elección del stack
tecnológico hasta la implementación de las reglas de negocio con las barajas.

Esta labor de documentación ha permitido realizar una reflexión crítica sobre el
trabajo realizado, asegurando que todo lo alcanzado en las fases anteriores quede
registrado. La memoria actúa como el soporte principal que demuestra la adquisición de
competencias en el uso de C\# y MongoDB, detallando cómo se ha resuelto el problema de
la variabilidad de datos mediante un sistema CRUD eficiente y escalable.

\section{Preparación de la presentación}

Un componente fundamental de esta etapa ha sido la preparación de la presentación
requerida para la asignatura. Este proceso ha consistido en sintetizar los aspectos
más relevantes del desarrollo de DeckBuilder para su exposición. Se han diseñado
materiales visuales que destacan los puntos fuertes de la aplicación, como la
integración de las APIs externas y la flexibilidad de la base de datos documental.

El objetivo de esta tarea es comunicar de forma clara y concisa los resultados
obtenidos y demostrar que el núcleo funcional desarrollado en la fase anterior es
operativo y responde a los objetivos planteados al inicio.

\section{Desarrollo del archivo README: Manual usuario y despliegue}

Como cierre técnico del proyecto, se ha elaborado un archivo README dentro del
repositorio del código fuente. Este documento se ha diseñado para cumplir una doble
función: servir como manual de usuario y como guía detallada de despliegue e
instalación. En él se explican paso a paso las instrucciones necesarias para que
cualquier usuario pueda configurar el entorno de ejecución, incluyendo la vinculación
con MongoDB, la instalación de dependencias de .NET 8 y el arranque del frontend.

Además de las instrucciones técnicas, el README incluye una guía de uso que describe
el flujo operativo de la aplicación y los enlaces de acceso al despliegue en Render.

\section{Resultados finales de la fase}

La conclusión de esta fase marca el fin del desarrollo de DeckBuilder. Con la memoria
técnica finalizada, la presentación preparada y el archivo README configurado, el
proyecto queda listo para su entrega definitiva. El resultado es una aplicación que
funciona de forma correcta, respaldada por una documentación para entender todo lo
realizado.
